% Bagian: computability
% Bab: computability-theory
% Subbagian: normal-form

\documentclass[../../../include/open-logic-section]{subfiles}

\begin{document}

\olfileid[id]{cmp}{thy}{nfm}
\olsection{Teorema Bentuk Normal}

Andaikan kita dapat mendeskripsikan definisi fungsi-fungsi dapat dihitung dan
menguji apakah suatu deskripsi yang diduga sebagai rekaman lengkap komputasi
fungsi tersebut pada suatu masukan memang benar. Maka, terlepas dari model
komputasi yang digunakan, sewajarnya kita dapat menentukan nilai setiap fungsi
dapat dihitung~$f$ pada setiap masukan~$x$ sebagai berikut:
\begin{enumerate}
  \item Telusuri semua deskripsi rekaman komputasi yang mungkin.
  \item Uji apakah suatu rekaman tertentu merupakan rekaman komputasi~$f(x)$.
  \item Jika rekaman yang benar telah ditemukan, ekstrak nilai~$f(x)$ darinya.
\end{enumerate}
Kebenaran prosedur ini merupakan isi teorema bentuk normal Kleene.

\begin{thm}[Teorema Bentuk Normal Kleene]
\ollabel{thm:normal-form}
Terdapat suatu relasi rekursif primitif~$T(e, x, s)$ dan suatu fungsi rekursif
primitif~$U(s)$ dengan sifat berikut: jika $f$ merupakan sembarang fungsi dapat
dihitung parsial, maka bagi suatu~$e$,
\[
f(x) \simeq U(\umin{s}{T(e, x, s)})
\]
untuk setiap~$x$.
\end{thm}

\begin{proof}[Sketsa Bukti]
Dalam setiap model komputasi, kita dapat secara ketat mendefinisikan suatu
deskripsi fungsi dapat dihitung~$f$ dan mengodekan deskripsi tersebut dengan
suatu bilangan asli~$e$. Kita juga dapat secara ketat mendefinisikan gagasan
``urutan komputasi'' yang merekam proses komputasi fungsi berindeks~$e$
untuk masukan~$x$. Urutan komputasi semacam itu juga dapat dikodekan sebagai
suatu bilangan~$s$. Hal ini dapat dilakukan sedemikian rupa sehingga
\begin{enumerate}
  \item relasi $T(e, x, s)$, yang berlaku jika dan hanya jika bilangan~$s$
  mengodekan urutan komputasi fungsi berindeks~$e$ pada masukan~$x$, dan
  \item fungsi $U(s)$, yang memetakan urutan komputasi berkode~$s$ ke hasil
  akhir komputasi tersebut,
\end{enumerate}
keduanya dapat dihitung. Bahkan, relasi~$T$ dan fungsi~$U$ bersifat rekursif
primitif.
\end{proof}

\begin{explain}
Untuk memberikan bukti ketat bagi Teorema Bentuk Normal, kita harus menetapkan
suatu model komputasi, melaksanakan pengodean deskripsi fungsi dapat dihitung
dan urutan komputasi secara terperinci, serta memverifikasi bahwa relasi~$T$
dan fungsi~$U$ bersifat rekursif primitif. Untuk kebanyakan penerapan, cukup
bahwa $T$ dan~$U$ dapat dihitung dan bahwa $U$~bersifat total.

Cara terbaik untuk mengingat bukti teorema bentuk normal barangkali melalui
slogan berikut: $\umin{s}{T(e, x, s)}$ mencari urutan komputasi fungsi
berindeks~$e$ pada masukan~$x$, dan $U$ mengembalikan keluaran urutan komputasi
tersebut jika urutan semacam itu dapat ditemukan.
\end{explain}

Jika model komputasinya berupa fungsi-fungsi rekursif parsial (seperti yang
semula digunakan Kleene), teorema tersebut menunjukkan bahwa hanya diperlukan
satu kali penggunaan pencarian tak berbatas, yakni satu operator
$\umin{y}{f(\vec x, y)}$, untuk mendefinisikan sembarang fungsi. Dalam
pengertian ini, teorema tersebut menunjukkan bahwa setiap fungsi rekursif
parsial mempunyai bentuk normal.

$T$ dan $U$ dapat digunakan untuk mendefinisikan enumerasi $\cfind{0}$,
$\cfind{1}$, $\cfind{2}$, \dots. Mulai sekarang, kita mengandaikan bahwa kita
telah menetapkan pilihan $T$ dan~$U$ yang sesuai, dan mengambil persamaan
\[
\cfind{e}(x) \simeq U(\umin{s}{T(e,x,s)})
\]
sebagai \emph{definisi} dari~$\cfind{e}$.

Berikut fakta berguna lainnya:

\begin{thm}
Setiap fungsi dapat dihitung parsial mempunyai tak berhingga banyak indeks.
\end{thm}

Sekali lagi, hal ini jelas secara intuitif. Dari sembarang (deskripsi suatu)
fungsi dapat dihitung, kita dapat membuat deskripsi lain yang menghitung fungsi
yang sama (pasangan masukan-keluarannya sama), tetapi melakukannya, misalnya,
dengan terlebih dahulu mengerjakan sesuatu yang tidak memengaruhi komputasi
(katakanlah menguji apakah $0 = 0$, menghitung sampai~$5$, dan sebagainya).
Indeks deskripsi yang diubah akan selalu berbeda dari indeks semula. Keduanya
merupakan indeks fungsi yang sama, hanya saja fungsi tersebut dihitung dengan
cara yang sedikit berbeda.

\end{document}
